\section{Why this benchmark exists, and what it found}
\label{sec:motivation}

A self-improvement benchmark reports a number. The number is only worth reading if the
measurement could have moved. This section states the four results that make the rest of the
report worth reading, each of which is about whether the instrument can see what it claims to
measure.

\subsection{The same runs, two scorers, opposite conclusions}
\label{sec:motivation:contrast}

The central result is a controlled comparison of two scoring rules on one experiment. Two cells
differ in a single environment variable and in nothing else --- same base model, same seed, same
sampling budget $k$, same task bank, same three arms, same code path. One scores each round with
the headroom-normalised best-of-$k$ rule the benchmark was designed around; the other scores the
identical generations by pass rate, the fraction of $k$ candidates that verify.

Across the ten rounds in which both the lineage and reset arms have reached correct-at-parity,
the headroom contrast reports a mean of $-0.0002$ and never departs from zero by more than
$0.004$. The same ten rounds report a mean of $+0.436$ under pass rate, with a minimum of
$+0.360$. The two ranges are separated by a factor of roughly ninety and do not approach one
another.

This is not a claim that pass rate is the correct metric; it has a ceiling of its own, reached in
three of four cells. It is a demonstration that a benchmark can return a confident null on a
difference its own runs contain. Table~\ref{tab:contrast} is the evidence, and it requires no
statistics to read.

\subsection{The blindness is scheduled, not accidental}

Best-of-$k$ is identical for a model that solves a task one time in $k$ and one that solves it
$k$ times in $k$. Self-training acts on exactly that quantity. The consequence is visible in the
arm levels rather than in any summary statistic: across all four headroom cells and both model
scales, the fourteen rounds in which the two arms sit within $0.02$ of each other average
$-0.0004$ and never exceed $0.004$, while the ten rounds with a gap average $+0.131$.

The contrast is therefore measuring the gap between the arms and nothing else, faithfully, at all
times. Best-of-$k$ compresses both arms onto correct-at-parity and holds them there, so the gap
it faithfully reports is zero. It recovers the moment a gap reopens: one cell reads $+0.002$ at
arm levels $0.50/0.50$ and $+0.201$ one round later at $0.50/0.30$.

The design consequence is the part worth carrying: a matched reset learner reaches parity within
three or four rounds, so the measurement's useful lifetime is fixed by the \emph{control} arm
rather than by the treatment. This cannot be repaired by choosing a better treatment or running
longer. A well-behaved control ends the measurement by behaving well.

\subsection{Three rungs did not produce nulls; they produced undefined measurements}
\label{sec:motivation:starved}

Three of the five rungs failed for one shared reason, and it is computable before a run rather
than diagnosed after. Each rung was built around a different channel, in a different lab, and in
each the channel received almost no input.

% Generated by scripts/make_kernel_results_tex.py. It was hand-typed until new rounds
% landed and moved it from 0.96/27 rounds/41% to 0.93/30/43% while the paper kept printing
% the old pair -- the drift this project regenerates everything else to avoid, on the number
% the report leads with.
% AUTO-GENERATED by scripts/make_kernel_results_tex.py -- do not edit by hand.
\begin{table}[H]\centering\small
\caption{The self-referential loop's throughput, measured. Each rung was designed to test whether a specific channel compounds; each channel was starved. The T2 figure is recomputed from the 30 rounds of the registered primary that are admissible under Amendments 5 and 6.}
\label{tab:starved}
\begin{tabular}{@{}llr@{}}
\toprule
Rung & What feeds the loop & Measured \\
\midrule
T2 (weight-RSI) & verified self-generated training examples & 0.93 per round, \textbf{zero in 43\%} of rounds \\
T3 (procedure-RSI) & kernels entering the archive after grading & \textbf{4} across six rounds, or \textbf{0} \\
T5 (self-play) & accepted model-authored tasks & \textbf{4 of 173} proposed (2.3\%) \\
\bottomrule
\end{tabular}
\end{table}


The arithmetic was available in advance. Expected training examples per round is
$n_{\mathrm{train}} \times k \times \text{yield}$. At the registered configuration this is
$3 \times 10 \times 0.032 = 0.96$, and the artifacts agree to two decimals (Table~\ref{tab:starved}
is recomputed at build time and moves as rounds land). We ran 57 trajectories and 228 rounds
of a loop that could not compound, and the calculation that says so is one line.

The claim this supports is not that recursive self-improvement fails. It is that \textbf{a
self-improvement benchmark must demonstrate that its loop has throughput before any contrast it
reports carries information}, and that this is a precondition to check rather than a result to
discover.

\subsection{A non-LLM baseline that scores below parity}

The most common objection to results in this domain is that a model reaching correct-at-parity is
simply failing. We answered it by submitting a non-LLM tool to the benchmark on its own terms.
\texttt{torch.compile(mode="max-autotune")}, scored exactly as a model candidate is, reaches a
median of $0.408$ --- below correct-at-parity --- because it is slower than default
\texttt{torch.compile} on 27 of 29 tasks. A model matching the compiled baseline is therefore
outperforming production autotuning on this bank, not failing to beat a weak baseline.

\subsection{What the reader should take from the rest of this report}

The four results above are about the instrument. The remainder of the report is the ladder that
produced them: five rungs, their designs, their controls, and every number with its uncertainty
and its scope. Section~\ref{sec:removed} lists every result that was removed or narrowed during
the work, with the evidence that forced each change, because a report that only shows its
surviving claims has hidden the part a reader most needs in order to trust the rest.
